% Generated by review_v53/build_joint_outcomes.py. Do not edit directly.
\begin{table}[!htb]
\centering\footnotesize
\setlength{\tabcolsep}{3pt}
\renewcommand{\arraystretch}{1.1}
\caption[Gemeinsame Ergebnisse zu Qualität, Durchsatz und Energieverbrauch]{Gemeinsame Ergebnisse für die sieben festgelegten ausgewählten „Split“ (S)- und „Full“ (F)-Bereitstellungen. Die Tabelle kombiniert die Verhältnisse von anhaltendem Durchsatz und Energieverbrauch mit direkten, paarweisen Qualitätsunterschieden im Vergleich zur dargestellten Full-Konfiguration.}
\label{tab:confirmation:joint_outcomes}
\begin{tabularx}{\linewidth}{@{}Y L{1.3cm}L{1.6cm}r r r@{}}
\toprule
Modell & Full & Split & \makecell[r]{$\Delta Q$ (pp)\\95\%-Paarintervall} & \makecell[r]{Durchsatz\\S/F} & \makecell[r]{Energie\\S/F} \\
\midrule
MobileNetV3-L & DX & DX / b135 & \makecell[r]{$+0.30$\\$[-0.04, +0.62]$} & 1.087 & 1.060 \\
RegNetX-1.6GF & DX & H10 / b123 & \makecell[r]{$+0.46$\\$[-0.10, +1.06]$} & 2.725 & 0.531 \\
ResNet50 & DX & H10 / b095 & \makecell[r]{$-0.38$\\$[-0.92, +0.12]$} & 1.561 & 1.115 \\
YOLO11l & DX & H10 / b120 & \makecell[r]{$-0.12$\\$[-0.63, +0.36]$} & 0.838 & 1.654 \\
YOLO26m & GPU (H8) & H10 / b043 & \makecell[r]{$-1.15$\\$[-1.48, -0.88]$} & 1.352 & 0.773 \\
YOLO26s & GPU (DX) & H8 / b021 & \makecell[r]{$-0.93$\\$[-1.21, -0.65]$} & 1.337 & 0.801 \\
YOLOv7 & DX & H8 / b066 & \makecell[r]{$-0.12$\\$[-0.31, +0.04]$} & 1.360 & 1.242 \\
\bottomrule
\end{tabularx}
\par\smallskip\begin{minipage}{\linewidth}\footnotesize $\Delta Q$ entspricht der Differenz zwischen „Split“ und „Full“ in Prozentpunkten: Top-1 für Klassifikatoren und COCO AP$_{50:95}$ für Detektoren. Die zugehörigen 95\%-Intervalle basieren auf 1\,000 gespeicherten Bootstrap-Stichproben aus denselben 5\,000 Bildern innerhalb jedes Paares. Tabelle~\ref{tab:confirmation:performance} enthält Durchsatzintervalle und Abkürzungen für den Einsatz. Tabelle~\ref{tab:confirmation:energy} spezifiziert die Energieaggregation und den Messumfang. Ein Durchsatz-S/F über eins und ein Energie-S/F unter eins sprechen für die Split-Aufteilung.\end{minipage}
\end{table}
